% Options for packages loaded elsewhere
\PassOptionsToPackage{unicode}{hyperref}
\PassOptionsToPackage{hyphens}{url}
%
\documentclass[
]{article}
\usepackage{amsmath,amssymb}
\usepackage{lmodern}
\usepackage{iftex}
\ifPDFTeX
  \usepackage[T1]{fontenc}
  \usepackage[utf8]{inputenc}
  \usepackage{textcomp} % provide euro and other symbols
\else % if luatex or xetex
  \usepackage{unicode-math}
  \defaultfontfeatures{Scale=MatchLowercase}
  \defaultfontfeatures[\rmfamily]{Ligatures=TeX,Scale=1}
\fi
% Use upquote if available, for straight quotes in verbatim environments
\IfFileExists{upquote.sty}{\usepackage{upquote}}{}
\IfFileExists{microtype.sty}{% use microtype if available
  \usepackage[]{microtype}
  \UseMicrotypeSet[protrusion]{basicmath} % disable protrusion for tt fonts
}{}
\makeatletter
\@ifundefined{KOMAClassName}{% if non-KOMA class
  \IfFileExists{parskip.sty}{%
    \usepackage{parskip}
  }{% else
    \setlength{\parindent}{0pt}
    \setlength{\parskip}{6pt plus 2pt minus 1pt}}
}{% if KOMA class
  \KOMAoptions{parskip=half}}
\makeatother
\usepackage{xcolor}
\IfFileExists{xurl.sty}{\usepackage{xurl}}{} % add URL line breaks if available
\IfFileExists{bookmark.sty}{\usepackage{bookmark}}{\usepackage{hyperref}}
\hypersetup{
  pdftitle={Homework 1: Multivariable Regression Analysis Report},
  pdfauthor={Xuange Liang},
  hidelinks,
  pdfcreator={LaTeX via pandoc}}
\urlstyle{same} % disable monospaced font for URLs
\usepackage{longtable,booktabs,array}
\usepackage{calc} % for calculating minipage widths
% Correct order of tables after \paragraph or \subparagraph
\usepackage{etoolbox}
\makeatletter
\patchcmd\longtable{\par}{\if@noskipsec\mbox{}\fi\par}{}{}
\makeatother
% Allow footnotes in longtable head/foot
\IfFileExists{footnotehyper.sty}{\usepackage{footnotehyper}}{\usepackage{footnote}}
\makesavenoteenv{longtable}
\setlength{\emergencystretch}{3em} % prevent overfull lines
\providecommand{\tightlist}{%
  \setlength{\itemsep}{0pt}\setlength{\parskip}{0pt}}
\setcounter{secnumdepth}{5}
% Shared LaTeX style for homework submissions
\usepackage{microtype}
\usepackage{booktabs}
\usepackage{longtable}
\usepackage{array}
\usepackage{caption}
\usepackage{fancyhdr}
\usepackage{titlesec}
\usepackage{setspace}
\usepackage{float}
\usepackage{graphicx}
\usepackage[table]{xcolor}
\usepackage{enumitem}
\usepackage[most]{tcolorbox}

\definecolor{TitleBlue}{HTML}{0B3C5D}
\definecolor{SoftGray}{HTML}{F6F8FA}

\setstretch{1.15}
\setlength{\parskip}{4pt}
\setlength{\parindent}{0pt}

\titleformat{\section}{\large\bfseries\color{TitleBlue}}{\thesection}{0.6em}{}
\titleformat{\subsection}{\normalsize\bfseries\color{TitleBlue}}{\thesubsection}{0.6em}{}

\captionsetup{
  font=small,
  labelfont=bf
}

\rowcolors{2}{SoftGray}{white}
\setlist[itemize]{leftmargin=*, topsep=2pt, itemsep=1pt}
\setlist[enumerate]{leftmargin=*, topsep=2pt, itemsep=1pt}

\pagestyle{fancy}
\fancyhf{}
\fancyhead[L]{\nouppercase{\leftmark}}
\fancyhead[R]{P8586 Applied Methods}
\fancyfoot[C]{\thepage}
\setlength{\headheight}{14pt}

\AtBeginDocument{
  \hypersetup{
    colorlinks=true,
    linkcolor=TitleBlue,
    urlcolor=TitleBlue,
    citecolor=TitleBlue
  }
}

% Keep title and TOC on same page, content starts on new page
\AtBeginDocument{
  \let\oldtableofcontents\tableofcontents
  \renewcommand{\tableofcontents}{%
    \oldtableofcontents
    \clearpage
  }
}
\ifLuaTeX
  \usepackage{selnolig}  % disable illegal ligatures
\fi

\title{Homework 1: Multivariable Regression Analysis}
\author{Xuange Liang}
\date{February 4, 2026}

\begin{document}
\maketitle

{
\setcounter{tocdepth}{3}
\tableofcontents
}
\hypertarget{research-question}{%
\section{Research Question}\label{research-question}}

\textbf{Objective}: To evaluate the association between intraventricular
hemorrhage (IVH) and neonatal mortality among very low birth weight
infants, adjusting for potential confounders.

\begin{itemize}
\tightlist
\item
  \textbf{Dataset}: \texttt{hw\_vlbw} (Very Low Birth Weight Infants,
  N=514)
\item
  \textbf{Exposure}: IVH (Intraventricular Hemorrhage)
\item
  \textbf{Outcome}: Death during hospitalization
\item
  \textbf{Study Design}: Retrospective cohort study
\end{itemize}

\begin{center}\rule{0.5\linewidth}{0.5pt}\end{center}

\hypertarget{methods}{%
\section{Methods}\label{methods}}

\hypertarget{data-source}{%
\subsection{1. Data Source}\label{data-source}}

\begin{itemize}
\tightlist
\item
  \textbf{Sample size}: 514 very low birth weight infants
\item
  \textbf{No missing data} on key variables
\item
  \textbf{Mortality rate}: 19.3\% (99/514)
\item
  \textbf{IVH prevalence}: 16.3\% (84/514)
\end{itemize}

\hypertarget{variable-definitions}{%
\subsection{2. Variable Definitions}\label{variable-definitions}}

\textbf{Primary Variables:} - Exposure: IVH (0=No, 1=Yes) - Outcome:
DEAD (0=Alive, 1=Dead)

\textbf{Covariates Considered:} - Pneumothorax (pneumo): 0=No, 1=Yes -
Delivery mode (delivery\_new): 1=C-Section, 2=Vaginal - Multiple
gestation (twin): 0=No, 1=Yes - Birth location (inout\_new): 1=Born at
Duke, 2=Transport - Birth weight (bwt): Continuous (grams) and
categorical - bwt\_cat: 1=\textless1000g, 2=1000-1299g, 3=1300-1599g -
Gestational age (gest): Continuous (weeks) and categorical - gest\_cat:
1=\textless28 weeks, 2=28-31 weeks, 3=≥32 weeks

\hypertarget{statistical-analysis}{%
\subsection{3. Statistical Analysis}\label{statistical-analysis}}

\textbf{Step 1: Descriptive Statistics} - Baseline characteristics
stratified by IVH status - Chi-square tests for categorical variables -
t-tests and Wilcoxon rank-sum tests for continuous variables

\textbf{Step 2: Unadjusted Association} - Logistic regression for odds
ratio - Poisson regression with log link for risk ratio - 2×2
contingency tables

\textbf{Step 3: Covariate Selection (Purposeful Selection Method)} -
Univariate screening: p \textless{} 0.25 threshold for candidate
variables (Hosmer, Lemeshow, \& Sturdivant) - Variables selected: IVH,
pneumo, inout\_new, bwt\_cat, gest\_cat - Variables excluded:
delivery\_new (p=0.40), twin (p=0.66)

\textbf{Step 4: Multivariable Model Building} - Initial model with all
candidate variables - Sequential removal of non-significant variables (p
\textgreater{} 0.10, per purposeful selection workflow) - Retention
criteria: removal changes IVH beta by \textless20\% (change-in-estimate
criterion for confounding assessment) - Variables removed: twin
(p=0.40), delivery\_new (p=0.62), inout\_new (p=0.93)

\textbf{Step 5: Final Model} - Adjusted for: pneumothorax, birth weight
category, gestational age category - Both logistic (OR) and Poisson (RR)
regression models estimated

\begin{center}\rule{0.5\linewidth}{0.5pt}\end{center}

\hypertarget{results}{%
\section{Results}\label{results}}

\hypertarget{table-1-baseline-characteristics-by-ivh-status}{%
\subsection{Table 1: Baseline Characteristics by IVH
Status}\label{table-1-baseline-characteristics-by-ivh-status}}

\begin{longtable}[]{@{}llll@{}}
\toprule
Characteristic & No IVH (n=430) & IVH (n=84) & p-value \\
\midrule
\endhead
\textbf{Demographics} & & & \\
Birth weight (g), mean±SD & 1101.4±246.0 & 964.1±266.7 &
\textless0.001 \\
Gestational age (wks), mean±SD & 29.1±2.3 & 27.4±2.2 & \textless0.001 \\
\textbf{Birth Weight Category, n (\%)} & & & 0.004 \\
\textless1000g & 146 (34.0\%) & 43 (51.2\%) & \\
1000-1299g & 170 (39.5\%) & 30 (35.7\%) & \\
1300-1599g & 114 (26.5\%) & 11 (13.1\%) & \\
\textbf{Gestational Age Category, n (\%)} & & & \textless0.001 \\
\textless28 weeks & 114 (26.5\%) & 40 (47.6\%) & \\
28-31 weeks & 255 (59.3\%) & 42 (50.0\%) & \\
≥32 weeks & 61 (14.2\%) & 2 (2.4\%) & \\
\textbf{Complications, n (\%)} & & & \\
Pneumothorax & 73 (17.0\%) & 34 (40.5\%) & \textless0.001 \\
Multiple gestation & 100 (23.3\%) & 7 (8.3\%) & 0.002 \\
\textbf{Delivery \& Location, n (\%)} & & & \\
C-Section delivery & 221 (51.4\%) & 27 (32.1\%) & 0.001 \\
Born at Duke & 366 (85.1\%) & 54 (64.3\%) & \textless0.001 \\
\textbf{Outcome} & & & \\
Death & 58 (13.5\%) & 41 (48.8\%) & \textless0.001 \\
\bottomrule
\end{longtable}

\textbf{Key Findings:} - Infants with IVH had significantly lower birth
weight (137g difference, p\textless0.001) - IVH was more common in
extremely premature infants (\textless28 weeks: 47.6\% vs 26.5\%) - IVH
co-occurred frequently with pneumothorax (40.5\% vs 17.0\%,
p\textless0.001) - Mortality was 3.6 times higher in IVH group (48.8\%
vs 13.5\%)

\begin{center}\rule{0.5\linewidth}{0.5pt}\end{center}

\hypertarget{table-2-mortality-by-baseline-characteristics}{%
\subsection{Table 2: Mortality by Baseline
Characteristics}\label{table-2-mortality-by-baseline-characteristics}}

\begin{longtable}[]{@{}lllll@{}}
\toprule
Characteristic & Alive (n=415) & Dead (n=99) & Mortality Rate &
p-value \\
\midrule
\endhead
\textbf{Exposure} & & & & \\
No IVH & 372 (89.6\%) & 58 (58.6\%) & 13.5\% & \textless0.001 \\
IVH & 43 (10.4\%) & 41 (41.4\%) & 48.8\% & \\
\textbf{Complications} & & & & \\
No pneumothorax & 359 (86.5\%) & 48 (48.5\%) & 11.8\% &
\textless0.001 \\
Pneumothorax & 56 (13.5\%) & 51 (51.5\%) & 47.7\% & \\
\textbf{Birth Weight} & & & & \textless0.001 \\
\textless1000g & 122 (29.4\%) & 67 (67.7\%) & 35.5\% & \\
1000-1299g & 173 (41.7\%) & 27 (27.3\%) & 13.5\% & \\
1300-1599g & 120 (28.9\%) & 5 (5.1\%) & 4.0\% & \\
\textbf{Gestational Age} & & & & \textless0.001 \\
\textless28 weeks & 93 (22.4\%) & 61 (61.6\%) & 39.6\% & \\
28-31 weeks & 261 (62.9\%) & 36 (36.4\%) & 12.1\% & \\
≥32 weeks & 61 (14.7\%) & 2 (2.0\%) & 3.2\% & \\
\bottomrule
\end{longtable}

\textbf{Interpretation of Table 2:} - Mortality was substantially higher
among infants with IVH than without IVH (48.8\% vs 13.5\%), showing a
strong crude exposure-outcome gradient. - Mortality was also much higher
among infants with pneumothorax (47.7\% vs 11.8\%), indicating severe
respiratory complications are strongly related to death. - A clear
dose-response pattern was present for birth weight: mortality decreased
from 35.5\% (\textless1000g) to 13.5\% (1000-1299g) to 4.0\%
(1300-1599g). - A similar gradient was observed by gestational age:
39.6\% (\textless28 weeks), 12.1\% (28-31 weeks), and 3.2\%
(\textgreater=32 weeks), supporting prematurity as a major mortality
driver. - Because IVH is associated with these strong mortality
predictors (Table 1), confounding is plausible and multivariable
adjustment is required.

\begin{center}\rule{0.5\linewidth}{0.5pt}\end{center}

\hypertarget{table-3-unadjusted-and-adjusted-association-between-ivh-and-mortality}{%
\subsection{Table 3: Unadjusted and Adjusted Association Between IVH and
Mortality}\label{table-3-unadjusted-and-adjusted-association-between-ivh-and-mortality}}

\begin{longtable}[]{@{}lllll@{}}
\toprule
Model & Measure & Estimate & 95\% CI & p-value \\
\midrule
\endhead
\textbf{Unadjusted} & & & & \\
Logistic & Odds Ratio & 6.116 & 3.674 - 10.179 & \textless0.001 \\
Poisson & Risk Ratio & 3.619 & 2.426 - 5.398 & \textless0.001 \\
\textbf{Adjusted¹} & & & & \\
Logistic & Odds Ratio & 4.149 & 2.291 - 7.516 & \textless0.001 \\
Poisson & Risk Ratio & 2.059 & 1.351 - 3.137 & 0.001 \\
\bottomrule
\end{longtable}

¹ Adjusted for pneumothorax, birth weight category, and gestational age
category

\textbf{Interpretation:} - After adjusting for key confounders, IVH
\textbf{increased the odds of death by 4.1-fold} (95\% CI: 2.3-7.5) - We
emphasize RR over OR for effect communication because death is not a
rare outcome in this cohort (99/514 = 19.3\%); with common outcomes, OR
can overstate the magnitude of association relative to RR. - Using the
more appropriate risk ratio for this context, IVH \textbf{doubled the
risk of death} (RR=2.1, 95\% CI: 1.4-3.1) - Adjustment attenuated the
effect by 32\% (OR: 6.1→4.1), indicating substantial but incomplete
confounding

\begin{center}\rule{0.5\linewidth}{0.5pt}\end{center}

\hypertarget{table-4-final-multivariable-model-results}{%
\subsection{Table 4: Final Multivariable Model
Results}\label{table-4-final-multivariable-model-results}}

\hypertarget{a.-logistic-regression-odds-ratios}{%
\subsubsection{A. Logistic Regression (Odds
Ratios)}\label{a.-logistic-regression-odds-ratios}}

\begin{longtable}[]{@{}llll@{}}
\toprule
Variable & Odds Ratio & 95\% CI & p-value \\
\midrule
\endhead
\textbf{IVH (Yes vs No)} & \textbf{4.149} & 2.291 - 7.516 &
\textless0.001 \\
Pneumothorax (Yes vs No) & 6.185 & 3.486 - 10.975 & \textless0.001 \\
\textbf{Birth Weight} & & & 0.0002 \\
\textless1000g vs 1300-1599g & 8.902 & 2.858 - 27.729 & 0.0002 \\
1000-1299g vs 1300-1599g & 3.216 & 1.113 - 9.293 & 0.031 \\
\textbf{Gestational Age} & & & 0.151 \\
\textless28 weeks vs ≥32 weeks & 3.691 & 0.749 - 18.177 & 0.108 \\
28-31 weeks vs ≥32 weeks & 2.239 & 0.484 - 10.352 & 0.302 \\
\bottomrule
\end{longtable}

\textbf{Model Performance:} - C-statistic: 0.853 (excellent
discrimination) - AIC: 373.1 - -2 Log Likelihood: 359.1

\hypertarget{b.-poisson-regression-risk-ratios}{%
\subsubsection{B. Poisson Regression (Risk
Ratios)}\label{b.-poisson-regression-risk-ratios}}

\begin{longtable}[]{@{}llll@{}}
\toprule
Variable & Risk Ratio & 95\% CI & p-value \\
\midrule
\endhead
\textbf{IVH (Yes vs No)} & \textbf{2.059} & 1.351 - 3.137 & 0.001 \\
Pneumothorax (Yes vs No) & 2.821 & 1.850 - 4.300 & \textless0.001 \\
\textbf{Birth Weight} & & & 0.004 \\
\textless1000g vs 1300-1599g & 4.861 & 1.779 - 13.283 & 0.002 \\
1000-1299g vs 1300-1599g & 2.692 & 1.025 - 7.069 & 0.044 \\
\textbf{Gestational Age} & & & 0.244 \\
\textless28 weeks vs ≥32 weeks & 3.143 & 0.710 - 13.923 & 0.132 \\
28-31 weeks vs ≥32 weeks & 2.333 & 0.555 - 9.809 & 0.248 \\
\bottomrule
\end{longtable}

\begin{center}\rule{0.5\linewidth}{0.5pt}\end{center}

\hypertarget{sensitivity-analysis}{%
\subsection{Sensitivity Analysis}\label{sensitivity-analysis}}

A simpler model including only IVH, pneumothorax, and birth weight
(without gestational age) was also tested:

\begin{longtable}[]{@{}llll@{}}
\toprule
Parameter & IVH OR (Full Model) & IVH OR (Simple Model) & \% Change \\
\midrule
\endhead
Beta coefficient & 1.4230 & 1.4750 & +3.7\% \\
Odds Ratio & 4.149 & 4.371 & +5.3\% \\
C-statistic & 0.853 & 0.842 & -1.3\% \\
\bottomrule
\end{longtable}

\textbf{Conclusion}: The IVH effect estimate is robust across model
specifications (\textless5\% change), supporting the validity of our
findings. The full model with gestational age is retained as the primary
analysis due to slightly better discrimination and theoretical
considerations.

\begin{center}\rule{0.5\linewidth}{0.5pt}\end{center}

\hypertarget{discussion}{%
\section{Discussion}\label{discussion}}

\hypertarget{main-findings}{%
\subsection{Main Findings}\label{main-findings}}

\begin{enumerate}
\def\labelenumi{\arabic{enumi}.}
\item
  \textbf{Strong Independent Effect}: After controlling for prematurity
  markers (birth weight, gestational age) and complications
  (pneumothorax), IVH remained strongly associated with neonatal
  mortality (adjusted OR=4.1, 95\% CI: 2.3-7.5).
\item
  \textbf{Substantial Confounding}: The crude association (OR=6.1) was
  attenuated by 32\% after adjustment, indicating that the relationship
  between IVH and mortality is partially explained by shared risk
  factors related to prematurity and illness severity.
\item
  \textbf{Pneumothorax as Key Predictor}: Pneumothorax showed the
  strongest association with mortality (OR=6.2), suggesting that
  respiratory complications are critical determinants of outcomes in
  this population.
\item
  \textbf{Dose-Response Relationship}: A clear gradient was observed for
  birth weight, with the lowest weight category (\textless1000g) having
  nearly 9-fold increased odds of death compared to the reference group
  (1300-1599g).
\item
  \textbf{Gestational Age Effect}: While gestational age categories
  showed non-significant associations in the multivariable model
  (p=0.15), this likely reflects collinearity with birth weight rather
  than lack of clinical importance.
\end{enumerate}

\hypertarget{why-rr-is-preferred-for-primary-interpretation}{%
\subsection{Why RR is Preferred for Primary
Interpretation}\label{why-rr-is-preferred-for-primary-interpretation}}

\begin{itemize}
\tightlist
\item
  The outcome (in-hospital death) occurred in 19.3\% of infants, which
  is above the usual ``rare outcome'' range where OR approximates RR.
\item
  OR is mathematically valid for logistic models but tends to exaggerate
  effect size when outcomes are common.
\item
  RR is more directly interpretable for cohort-style risk communication
  (``times the risk'') and better aligns with clinical/public-health
  interpretation.
\item
  Therefore, OR and RR are both reported for completeness, but RR is
  used as the primary effect measure in interpretation.
\end{itemize}

\hypertarget{clinical-implications}{%
\subsection{Clinical Implications}\label{clinical-implications}}

\begin{itemize}
\tightlist
\item
  IVH represents a \textbf{major independent risk factor} for mortality
  in very low birth weight infants
\item
  Early detection and management of IVH should be prioritized
\item
  Co-occurrence of IVH and pneumothorax identifies an extremely
  high-risk subgroup (likely \textgreater50\% mortality)
\item
  Prevention strategies targeting the most premature infants
  (\textless28 weeks, \textless1000g) may have the greatest impact on
  reducing mortality
\end{itemize}

\hypertarget{strengths}{%
\subsection{Strengths}\label{strengths}}

\begin{itemize}
\tightlist
\item
  Complete data with no missing values
\item
  Appropriate statistical methods (purposeful selection, both OR and RR
  reported)
\item
  Robust effect estimates across model specifications
\item
  Excellent model discrimination (c-statistic=0.853)
\end{itemize}

\hypertarget{limitations}{%
\subsection{Limitations}\label{limitations}}

\begin{enumerate}
\def\labelenumi{\arabic{enumi}.}
\tightlist
\item
  \textbf{Collinearity}: High correlation between birth weight and
  gestational age limits ability to disentangle their independent
  effects
\item
  \textbf{Residual Confounding}: Unmeasured factors (e.g., severity of
  IVH, ventilator settings) may still confound the association
\item
  \textbf{Sample Size}: Wide confidence intervals for some estimates,
  particularly gestational age categories
\item
  \textbf{Temporal Ambiguity}: Cannot determine if some complications
  (e.g., pneumothorax) occurred before or after IVH
\end{enumerate}

\begin{center}\rule{0.5\linewidth}{0.5pt}\end{center}

\hypertarget{conclusions}{%
\section{Conclusions}\label{conclusions}}

Intraventricular hemorrhage is a strong, independent predictor of
neonatal mortality among very low birth weight infants. After adjusting
for pneumothorax, birth weight, and gestational age, infants with IVH
had \textbf{4.1 times higher odds} (OR=4.1, 95\% CI: 2.3-7.5) and
\textbf{2.1 times higher risk} (RR=2.1, 95\% CI: 1.4-3.1) of death
compared to those without IVH. These findings underscore the critical
importance of preventing and managing IVH in this vulnerable population.

\begin{center}\rule{0.5\linewidth}{0.5pt}\end{center}

\hypertarget{references}{%
\section{References}\label{references}}

\begin{enumerate}
\def\labelenumi{\arabic{enumi}.}
\tightlist
\item
  Hosmer DW, Lemeshow S, Sturdivant RX. \emph{Applied Logistic
  Regression}. 3rd ed.~New York: Wiley; 2013. (Purposeful selection
  workflow, including p\textless0.25 screening and iterative model
  reduction.)
\item
  Mickey RM, Greenland S. The impact of confounder selection criteria on
  effect estimation. \emph{Am J Epidemiol}. 1989;129(1):125-137.
  (Rationale for change-in-estimate confounder assessment.)
\item
  Zhang J, Yu KF. What's the relative risk? A method of correcting the
  odds ratio in cohort studies of common outcomes. \emph{JAMA}.
  1998;280(19):1690-1691. (Why OR diverges from RR when outcome is
  common.)
\item
  Zou G. A modified Poisson regression approach to prospective studies
  with binary data. \emph{Am J Epidemiol}. 2004;159(7):702-706. (RR
  estimation using Poisson model with log link for binary outcomes.)
\end{enumerate}

\textbf{Software:} - SAS 9.4 (SAS Institute, Cary, NC) - Procedures:
PROC LOGISTIC, PROC GENMOD, PROC FREQ, PROC MEANS

\begin{center}\rule{0.5\linewidth}{0.5pt}\end{center}

\hypertarget{appendix-variable-coding}{%
\section{Appendix: Variable Coding}\label{appendix-variable-coding}}

\begin{verbatim}
IVH:          0 = No IVH, 1 = IVH present
DEAD:         0 = Alive at discharge, 1 = Dead
pneumo:       0 = No pneumothorax, 1 = Pneumothorax
bwt_cat:      1 = <1000g, 2 = 1000-1299g, 3 = 1300-1599g (reference)
gest_cat:     1 = <28 weeks, 2 = 28-31 weeks, 3 = ≥32 weeks (reference)
delivery_new: 1 = C-Section (reference), 2 = Vaginal
twin:         0 = Singleton (reference), 1 = Multiple gestation
inout_new:    1 = Born at Duke (reference), 2 = Transport
\end{verbatim}

\begin{center}\rule{0.5\linewidth}{0.5pt}\end{center}

\textbf{Report Generated}: February 4, 2026\\
\textbf{Analysis Code}: \url{HW1_Analysis.sas}

\end{document}
